This study develops and evaluates a hybrid monitoring scheme for nonlinear, serially dependent multivariate processes that combines a
one-dimensional convolutional autoencoder (CAE), a one-class support vector machine (OC-SVM) and
a multivariate exponentially weighted moving average (MEWMA) control chart. For every window of
32 observations, the CAE's reconstruction error and the OC-SVM's anomaly score on the CAE
features form a two-dimensional monitoring vector that the MEWMA accumulates over time. The
scheme is studied on a simulated nonlinear process driven by a latent vector autoregression,
using five independent datasets for representation learning, detector fitting, empirical
calibration of the control limit, independent in-control (IC) validation, and out-of-control (OOC)
evaluation. The results reveal that the hybrid chart has a shorter OOC average run-length (\ARL{}) than both of its single-component ablations at
every shift, but it does not outperform the classical MEWMA: the two are indistinguishable at the smallest
shift, and the classical chart detects faster at every shift of 0.4 or more. Post-hoc
diagnostics attribute this mainly to the weak response of distance-type learned features when the
mean shifts in a fixed direction, amplified at large shifts by the startup behaviour of the charts. Simulated
Extended Tennessee Eastman Process data were used to demonstrate the application of the proposed
method. The frozen design was refitted to 31 process variables and both charts were
calibrated to a 5\% probability of at least one false alarm over 568 nominal decisions instead of an IC \ARL{} of 370. The classical MEWMA again detected the feed-composition fault IDV(2) sooner at every severity, by 4.8 to
10.1 monitoring decisions on average in onset-referenced conditional detection delay, while the
hybrid chart's median delay was 17 to 39 decisions (0.80 to 1.90 hours after onset), depending on
severity.

\noindent\textbf{Keywords:} Statistical process control; MEWMA chart; Convolutional autoencoder;
One-class support vector machine; Multivariate process monitoring; Anomaly detection;
Tennessee Eastman Process.